\documentclass[a4paper,11pt,twoside]{amsart}
\usepackage[margin=1in]{geometry}
\usepackage{mathtools}
\usepackage{amsfonts}
\usepackage{amsthm}
\usepackage{amsmath}
\usepackage{todonotes}
\usepackage{enumerate}
\usepackage{graphicx}
\usepackage{subcaption}
\usepackage{url}
\usepackage{tabularx}
\usepackage{biblatex}

\let\originalleft\left
\let\originalright\right
\let\vec\mathbf
\renewcommand{\left}{\mathopen{}\mathclose\bgroup\originalleft}
\renewcommand{\right}{\aftergroup\egroup\originalright}
\renewcommand{\Re}{\operatorname{Re}}
\renewcommand{\Im}{\operatorname{Im}}
\newcommand{\Log}{\operatorname{Log}}
\newcommand{\Arg}{\operatorname{Arg}}
\newcommand{\C}{\mathbb{C}}
\newcommand{\R}{\mathbb{R}}
\newcommand{\N}{\mathbb{N}}
\newcommand{\eps}{\varepsilon}

\newcommand\numberthis{\addtocounter{equation}{1}\tag{\theequation}}

\newtheorem*{theorem}{Theorem}
\newtheorem*{corollary}{Corollary}
\newtheorem{lemma}{Lemma}
\newtheorem*{fact}{Fact}
\theoremstyle{definition}
\newtheorem*{definition}{Definition}

%\setlength{\oddsidemargin}{.5in}
%\setlength{\evensidemargin}{.5in}
%\setlength{\textwidth}{\paperwidth}
%\addtolength{\textwidth}{-1in}
%\setlength\evensidemargin\oddsidemargin

\addbibresource{references.bib}

\begin{document}
\title{Proving the spectral theorem}
\author{Alex Dobner}
\address{Department of Mathematics, University of Michigan \\
530 Church St\\
Ann Arbor, MI 48109}
\email{adobner@umich.edu}
\maketitle

In this note I'll summarize (for myself) how to think about and prove the spectral theorem in various guises.

The most basic version---as one learns in a linear algebra class---covers how a self-adjoint operator acts upon a finite dimensional complex inner product space (or how a self-adjoint matrix acts on $\C^n$ if you like).
\begin{theorem}
Let $T \colon V \to V$ be a self-adjoint operator on a finite dimensional inner product space. There exists an orthonormal eigenbasis for $T$.
\end{theorem}

A slightly different way to say this is that one can decompose $V$ into a direct sum of orthogonal components (the eigenspaces), and on each of these components $T$ simply acts like a multiplication by a scalar (the eigenvalue of that eigenspace).

\begin{proof}[Proof sketch]
One starts by locating a single eigenvalue of $T$. This can easily be done by considering the characteristic polynomial of $T$. The associated eigenspace is clearly $T$-invariant, so the orthogonal complement of this eigenspace is also $T$ invariant (this is where we use the self-adjointness assumption). Now we restrict $T$ to this invariant subspace which is of smaller dimension and repeat. Of course if at any point the orthogonal complement is empty then we're done.
\end{proof}

\textbf{Note:} The same statement holds if we assume $T$ is unitary instead, and the proof is almost exactly the same because the unitary assumption also implies that the orthogonal complement of an eigenspace is $T$-invariant. One can also generalize both of these statements to normal operators, although this requires a small additional trick.\footnote{The trick: one can write any normal operator as the sum of two commuting self-adjoint operators.}

\section{Generalizing to compact operators}
Suppose we want to generalize this statement to the infinite dimensional spaces. The natural way to do this is to now assume the space is a Hilbert space.\footnote{The completeness assumption for Hilbert spaces is responsible for making lots of things that are true in the finite dimensional setting carry over to the infinite dimensional setting. One could also consider operators on Banach spaces but this is less common.} The notions of self-adjoint/unitary/normal operators all carry over to this setting.

Note that in the proof of the finite dimensional spectral theorem above we used the finiteness in two ways: (1) it allowed us to use the characteristic polynomial to get an eigenvalue, and (2) it guaranteed that our procedure would terminate eventually. Here is the simplest generalization to the infinite dimensional setting where we can get around these issues in a somewhat natural way.
\begin{theorem}
Let $T \colon V \to V$ be a compact self-adjoint operator on a Hilbert space $V$. Then $T$ has an orthonormal eigenbasis. Moreover, the eigenspaces are finite dimensional (except for possibly the $\lambda=0$ eigenspace). Also, if the sequence of $\lambda_n$'s is infinite then it converges to zero.
\end{theorem}


\begin{proof}[Proof sketch]
One strategy is as follows: show that (1) if $\lambda$ is a nonzero element of spectrum of a compact operator, then it is an eigenvalue, and that (2) the spectral radius of a self-adjoint operator equals the operator norm. Combining these two facts, if $T$ is a nonzero compact self-adjoint operator then it has an eigenvalue $\lambda_1$ such that $||T||=|\lambda_1|.$ The orthogonal complement of the eigenspace associated to $\lambda_1$ is $T$-invariant by the self-adjointness, so we may restrict $T$ to this space and repeat. (Note that the eigenspace must be finite-dimensional otherwise this would violate the compactness of $T$.)

There are now two possibilities for how this process goes:
\begin{enumerate}
    \item If $T$ ends up being the zero operator at some point, then we stop and pick an orthonormal basis for this (infinite dimensional) space.
    \item If $T$ is never the zero operator then it must be the case that $\lambda_n \to 0$, otherwise the compactness of $T$ fails. Now note that if $x \in V$ is orthogonal to all the eigenspaces selected in this process, then we must have that $||Tx||\leq |\lambda_n| ||x|| \to 0$, so $x$ is in the 0-eigenspace.
\end{enumerate}

\end{proof}

Note that the compactness assumption allowed us to carry out a proof which was rather similar to the finite dimensional case. Unfortunately, without compactness things get much worse and a completely different approach seems to be needed.

\section{Alternative Method}
\subsection{A bird's-eye view}
What is the idea behind the spectral theorem? It says that given some operator $T\colon V \to V$ which is sufficiently nice we can ``chop up'' $V$ into orthogonal subspaces, such that on each subspace $T$ acts like multiplication by a scalar.

If $v \in V$ is not in a single eigenspace then it is some linear combination of eigenvectors (in the finite/compact case at least), and so we can think of $T(v)$ as a certain average of the different multiplication operators.

If we remove the compactness assumption, it turns out that we can no longer hope to chop up the space $V$ into `discrete' components. However, the action of $T$ on some $v \in V$ might still be able to viewed as a kind of `continuous' average of multiplication operators on the different `components' of $v$. (In other words, $T$ may have ``continuous spectrum''.)

\subsection{Baby functional calculus}
In order to discover how to prove more general versions of the spectral theorem we first observe a certain consequence of the versions we have already seen known as the \emph{functional calculus}.

\begin{theorem}[Baby functional calculus]
Let $T \colon V \to V$ be a self-adjoint operator on a finite dimensional complex vector space. Let $\sigma(T)$ denote the set of eigenvalues of $T$. There is a unique *-homomorphism $\Phi \colon C(\sigma(T)) \to \mathcal{B}(V)$ such that $\Phi(\mathrm{id}) = T.$
\begin{proof}
The existence is easy to prove using the spectral theorem. Write $T = U \mathrm{diag}(\lambda_1, \ldots,\lambda_n) U^{-1}$ where $U$ is unitary. Then let $\Phi(f) = U \mathrm{diag}(f(\lambda_1), \ldots,f(\lambda_n)) U^{-1}.$

One method to show uniqueness is to show that for any eigenvalue $\lambda$, the operator $\Phi(1_{\{\lambda\}})$ must be the orthogonal projection onto the $\lambda$ eigenspace. This then uniquely determines $\Phi.$
\end{proof}
\end{theorem}

In the proof above we have deduced the functional calculus from the spectral theorem, but it turns out that it is also possible to deduce the spectral theorem from the functional calculus. How do we do this? Well, in this ``baby version'' where $V$ is finite dimensional, there's not too much technical difficulty in doing so because $\sigma(T)$ is a finite set, so if we're handed the function $\Phi$ we know that its behavior is completely determined by how it acts on the characteristic functions $1_{\{\lambda\}}$ where $\lambda \in \sigma(T)$.

Indeed, given a $T \colon V \to V$ as above, note that the identity function on $\sigma(T)$ can be written as $\sum_{\lambda \in \sigma(T)} \lambda 1_{\{\lambda\}}.$ So
\[
T = \Phi(\mathrm{id}) = \Phi\left(\sum_{\lambda \in \sigma(T)} \lambda 1_{\{\lambda\}}\right) =
\sum_{\lambda \in \sigma(T)} \lambda \Phi(1_{\{\lambda\}}).
\]
It's now easy to check that the $\Phi(1_{\{\lambda\}})$ operators are self-adjoint projections whose images are the eigenspaces of $T$.

\subsection{Generalizing}

What we've seen above now gives an alternative method of proving the spectral theorem.
\begin{enumerate}
    \item Given $T \colon V \to V$, construct a functional calculus $\Phi \colon C(\sigma(T)) \to \mathcal{B}(V)$ which is a *-homomorphism such that $\Phi(\mathrm{id})=T$.
    \item Express the map $\Phi$ ``concretely'' in terms of projections onto eigenspaces.
\end{enumerate}
Note that for operators on infinite dimensional spaces, the spectrum $\sigma(T)$ may no longer be a discrete set of points, so the second step will not be as simple as before because it won't be the case that $C(\sigma(T))$ is spanned by a finite collection of characteristic functions. In this scenario the ``concrete'' expression of $\Phi$ will not be as a finite sum of projection but as an integral!

Interestingly, now that we've isolated our two steps, we can step back and try to turn both of them into more general statements. Essentially we are asking then: what are the key properties that we are requiring to prove the statements like these?

\subsection{Abstract functional calculus}
Here is a much more abstract version of step 1.

TODO: Give a definition of C*-algebra

\begin{theorem}
Let $a$ be a normal element of a C*-algebra $\mathcal{A}$. Then there is a unique *-homomorphism $\Phi \colon C(\sigma(a)) \to C^*(1,a) \subset \mathcal{A}$ such that $\Phi(\mathrm{id}) = a.$ Moreover, $\Phi$ is an isomorphism of C*-algebras.
\end{theorem}

If $\mathcal{A}$ is the C*-algebra of bounded operators on a Hilbert space then this is exactly the statement we need.

\section{Abstract Riesz Representation Theorem}
Here's a version of the Riesz representation theorem for *-representations.
\begin{theorem}
    Let $X$ be a compact hausdorff space and let $\Omega$ denote the Borel $\sigma$-algebra of $X$. For any *-representation, $\pi \colon C(X) \to \mathcal{B}(V)$, there exists a unique regular spectral measure $E \colon \Omega \to \mathcal{B}(V)$ such that
    \[
    \pi(f) = \int f \, dE.
    \]
\end{theorem}

The proofs I've found use a ``weak'' approach where, given $x,y \in V$, one considers the linear functional $f \mapsto \langle \pi(f)x,y\rangle$ on functions $f \in C(X)$. This functional yields a complex measure $E_{x,y}$ by the Riesz representation theorem. Now, given a bounded, \emph{measurable} function $g$, we can consider the map $(x,y) \mapsto \int g E_{x,y}$. By the other Riesz representation theorem, this map is induced by an element of $\mathcal{B}(V)$ which is what we assert as the correct definition for $\int g\,  dE.$

\section{Von Neumann Algebras}

TODO: Give a definition of von Neumann algebra as a C* algebra which is closed under supremums etc. \cite{MR861016}

\nocite{*}
\printbibliography
\end{document}
